\documentclass[10pt,a4paper]{amsart}
\usepackage[margin=19mm]{geometry}
\usepackage{amsmath,amssymb,mathtools}
\usepackage[scr=rsfs,cal=euler]{mathalfa}
\usepackage{xfrac}
\usepackage[unicode,hidelinks,bookmarks=false]{hyperref}
\newcommand{\ie}{i.e.\ }
\newcommand{\cf}{cf.\ }
\newcommand{\resp}{respectively\ }
\newcommand{\Subsection}{\subsection}
\providecommand{\nobreakdash}{\nobreak\hbox{-}}
\setlength{\emergencystretch}{2em}
\multlinegap0pt
\usepackage{amssymb}
\usepackage{xcolor}
\usepackage{enumerate}
\usepackage{tikz}
\usepackage{float}
\newtheorem{lemma}{Lemma}
\newtheorem{proposition}{Proposition}
\usepackage{cleveref}
\crefname{lemma}{Lemma}{Lemmas}
\crefname{proposition}{Proposition}{Propositions}
\title{On probabilistic identities and \\ coset identities in pro-$p$ groups}
\author{Steffen Kionke, Nowras Otmen, Tommaso Toti, Matteo Vannacci, Thomas Weigel}
\date{Source: arXiv:2507.19086v2 (CC BY 4.0)}
\begin{document}
\maketitle

\noindent\emph{Source and licence.} On probabilistic identities and \\ coset identities in pro-$p$ groups. Version arXiv:2507.19086v2 (CC BY 4.0), distributed by arXiv under the Creative Commons Attribution 4.0 International licence, http://creativecommons.org/licenses/by/4.0/, as recorded in the arXiv licence metadata for this exact version and shown on the abstract page https://arxiv.org/abs/2507.19086v2. Full licence notice: \texttt{licenses/arxiv-2507.19086-CC-BY-4.0.txt}.\par\smallskip

\noindent\textit{Editorial scope.} Counted unit: the article's proposition pr:ccmeasure (source0.tex lines 606--616). The article's complete original proof of it is delivered with it (source0.tex lines 617--636, one \texttt{proof} environment of 20 lines). No mathematics is written, repaired or re-derived: the statement and the proof are the article's own lines, in the article's own notation, and no retained line is substituted or re-flowed.  Retained uncounted as the source-local context the proof needs: lines 248--257.  Nothing was omitted from any retained span, so the package carries no omission record.  No retained line contains a citation, so the wrapper carries no bibliography and no bibliography entry.  No retained line contains a figure environment, an image-inclusion command or drawing code, so no image is delivered and the package declares no asset dependency.  Editorial formatting only, in addition: an AMS \texttt{amsart} preamble that restates the article's own declarations and macros used by the retained lines, namely the environments \texttt{lemma}, \texttt{proposition} on separate counters, together with the standard packages those lines need.  The retained environments print in this document as Lemma 1, Proposition 1 in order of appearance, whereas the article prints the same items with the numbering of its own full document.  The complete native source of the article as distributed in the arXiv e-print for this exact version is bundled unchanged as \texttt{sources/182188/source0.tex}.
\begin{lemma}\label{quotient}
	   	Let $G$ be a profinite group and let $X$ be a closed subset of $G$. If $K\unlhd G$ and $\pi\colon G\to G/K$ is the natural projection, then \begin{equation}\label{eq:leqquotient}
	   		\mu_G(X)\leq\mu_{G/K}(\pi(X)).
	   	\end{equation}
	   	 If moreover $G$ is countably based, then 
         \begin{equation}\label{eq:measurefinitequotients}
     	\mu_G(X)=\underset{N\unlhd_oG}{\mathrm{inf}}\frac{|XN/N|}{|G/N|}=\underset{i\geq 1}{\mathrm{inf}}\frac{|XN_i/N_i|}{|G/N_i|},
     \end{equation}
     for every $\{N_i\}_{i\geq 1}$ descending chain of open normal subgroups of $G$ with trivial intersection.
\end{lemma}

	\begin{proposition}\label{pr:ccmeasure}
		Let $G$ be a countably based profinite group and let $x\in G$. Given $\{N_i\}_{i\geq 1}$ a descending chain of open normal subgroups of $G$ with trivial intersection, define 
		\begin{align*}
			& c_1(x)=\underset{i\geq 1}{\mathrm{sup}}|C_{G/N_i}(xN_i)|;\\
			& c_2(x)=\underset{N\unlhd_oG}{\mathrm{sup}}|C_{G/N}(xN)|;\\
			& c_3(x)=\underset{K\unlhd G}{\mathrm{sup}}|C_{G/K}(xK)|,
		\end{align*}
		where $c_3(x)=\infty$ if $C_{G/K}(xK)$ is infinite for some $K\unlhd G$. Then $c_1(x)=c_2(x)=c_3(x):=c(x)$ and \begin{equation}\label{eq:ccmeasure}
			\mu(x^G)=\frac{1}{c(x)}.
		\end{equation}
	\end{proposition}

	\begin{proof}
		Since for any $N\unlhd_o G$ one has $x^GN/N=(xN)^{G/N}$, it follows that
		\begin{equation*}\label{eq:orbstab}
			|x^GN/N|=|(xN)^{G/N}|=|G/N:C_{G/N}(xN)|=\frac{|G/N|}{|C_{G/N}(xN)|}. 
		\end{equation*} 
		By \cref{quotient}, one obtains that 
        
		\begin{equation*}
			\mu(x^G)=\underset{i\geq 1}{\mathrm{inf}}\frac{1}{|C_{G/N_i}(xN_i)|}=\underset{N\unlhd_oG}{\mathrm{inf}}\frac{1}{|C_{G/N}(xN)|}. 
		\end{equation*} Therefore $c_1(x)=c_2(x)$ and $\mu(x^G)=1/c_i(x)$ for $i=1,2$. Since $c_2(x)\leq c_3(x)$, it remains to verify that $c_3(x)\leq c_2(x)$. Let $K\unlhd G$ and consider the quotient group $G/K\simeq \underset{N\unlhd_oG}\varprojlim\frac{G}{KN}$.
        Since $C_{G/K}(xK)$ is a closed subgroup of $G/K$, one has that 
		\begin{equation}\label{eq:centrquotient}
			C_{G/K}(xK)\simeq\underset{N\unlhd_oG}\varprojlim\frac{C_{G/K}(xK)KN}{KN}.
		\end{equation}
		Moreover, $C_{G/K}(xK)KN/KN\leq C_{G/KN}(xKN)$ for every $N\unlhd_oG$. Therefore \cref{eq:centrquotient} implies the following chain of inequalities 
		\begin{equation*}
			|C_{G/K}(xK)|=\underset{N\unlhd_oG}{\mathrm{sup}}\left\vert\frac{C_{G/K}(xK)KN}{KN}\right\vert\leq \underset{N\unlhd_oG}{\mathrm{sup}}|C_{G/KN}(xKN)|\leq c_2(x).
		\end{equation*}
		Since the previous inequality holds for every $K\unlhd G$, it follows that $c_3(x)\leq c_2(x)$.
	\end{proof}

\end{document}
